% --- Archon physics formalization source begin ---
% archon:physics
% archon:covers PhyXMiniProblems/problem_phyx_mini_0046.lean
% archon:source-report reports/phyx_mini/problem_phyx_mini_0046.source.json

\chapter{Physics problem phyx\_mini\_0046}
\label{ch:PhyXMiniProblems_problem_phyx_mini_0046}

\paragraph{Problem source.}
\#\# Physical scenario

The figure shows a two-slit interference experiment in which the slits are 0.200\textbackslash{} \textbackslash{}mathrm\{mm\} apart and the screen is 1.00\textbackslash{} \textbackslash{}mathrm\{m\} from the slits. The \textbackslash{}( m = 3 \textbackslash{}) bright fringe in the figure is 9.49\textbackslash{} \textbackslash{}mathrm\{mm\} from the central bright fringe.

\#\# Question

Find the wavelength of the light.

\#\# Answer choices

- A: A: 643nm

- B: B: 633nm

- C: C: 639nm

- D: D: 533nm

\#\# Figure caption (auxiliary; use the image as primary evidence)

The image illustrates a double-slit interference setup.

1. **Components and Arrangement**:
   - A light source on the left emits a beam toward a barrier with a single slit.
   - The light then hits a second barrier containing two slits, labeled as "Slits."
   - The distance between the slits, \textbackslash{}( d \textbackslash{}), is denoted as 0.200 mm.
   - A screen is positioned on the right to capture the interference pattern, with the distance from the slits to the screen \textbackslash{}( R \textbackslash{}) being 1.00 m.

2. **Interference Pattern**:
   - The light passes through the slits, creating a pattern on the screen.
   - The pattern is shown as a series of bright and dark fringes.
   - The spacing between adjacent fringes on the screen is labeled as 9.49 mm.

3. **Coordinates and Labels**:
   - The x-axis runs horizontally, while the y-axis is vertical through the screen.
   - Fringe orders are labeled on the screen, with \textbackslash{}( m \textbackslash{}) values from -3 to +3, indicating the order of the bright fringes.

This setup demonstrates the classic double-slit experiment, showing how light exhibits wave-like interference.

\#\# Dataset metadata

Wave Optics; Physical Model Grounding Reasoning, Multi-Formula Reasoning

\paragraph{Recorded answer/context.}
B: B: 633nm

\paragraph{Figure/image path.}
/root/proposal\_for\_physic/science-mango/phyx\_mini\_run/phyx\_data/test\_image/46.png

\paragraph{Formalization target.}
create a compiling Lean file with sorry bodies at `PhyXMiniProblems/problem\_phyx\_mini\_0046.lean`. The Lean declarations must preserve the physical quantities, dimensions or dimensional roles, figure labels, governing-law hypotheses, and final relation expressed by this problem.
Use Mathlib/Physlib names found through LeanExplore where available. If a physics API is missing, introduce faithful local abstractions rather than scalar placeholder aliases.

\begin{theorem}[Physics formalization target]
\label{thm:physics:phyx_mini_0046:target}
\lean{PhyXMiniProblems.ProblemPhyXMini0046.wavelength_rounds_to_recordedAnswerB}
\uses{def:physics:phyx-mini-0046:phyxminiproblems-problemphyxmini0046-nanometersvalue, def:physics:phyx-mini-0046:phyxminiproblems-problemphyxmini0046-doubleslitapparatus, def:physics:phyx-mini-0046:phyxminiproblems-problemphyxmini0046-brightfringeobservation, def:physics:phyx-mini-0046:phyxminiproblems-problemphyxmini0046-hasstatedfigurereadouts, def:physics:phyx-mini-0046:phyxminiproblems-problemphyxmini0046-satisfiesdoubleslitbrightfringelaws, def:physics:phyx-mini-0046:phyxminiproblems-problemphyxmini0046-answerwavelengthnanometers}
For slit separation \(0.200\,\mathrm{mm}\), screen distance
\(1.00\,\mathrm m\), and the third bright fringe at \(9.49\,\mathrm{mm}\),
the exact screen geometry and double-slit maximum law give a wavelength whose
nanometre readout rounds to \(633\), answer B.
\end{theorem}
\begin{proof}
Let \(\theta\) be the acute fringe angle.  The screen relation gives
\(\tan\theta=0.00949\), hence
\(\sin\theta=0.00949/\sqrt{1+0.00949^2}\).  The \(m=3\) bright-fringe law
then gives
\[
 \lambda=\frac{0.0002}{3}\,
   \frac{0.00949}{\sqrt{1+0.00949^2}}\ {\rm m}.
\]
All quantities are positive.  Square the equivalent nonnegative inequalities
and normalize the rational expressions to show that the nanometre readout
lies in the rounding interval \([632.5,633.5)\).  Thus its nearest integer is
the value printed for B.
\end{proof}

% --- Archon declaration topology begin ---
\section{Declaration topology}
\begin{definition}[Dim Length]
  \label{def:physics:phyx-mini-0046:phyxminiproblems-problemphyxmini0046-dimlength}
  \lean{PhyXMiniProblems.ProblemPhyXMini0046.DimLength}
  A signed physical length, represented independently of a chosen unit system
\end{definition}
\begin{proof}
  This is the declared model definition of Dim Length.
\end{proof}
\begin{definition}[length Value In]
  \label{def:physics:phyx-mini-0046:phyxminiproblems-problemphyxmini0046-lengthvaluein}
  \lean{PhyXMiniProblems.ProblemPhyXMini0046.lengthValueIn}
  \uses{def:physics:phyx-mini-0046:phyxminiproblems-problemphyxmini0046-dimlength}
  Read a physical length as a real scalar in the specified length unit
\end{definition}
\begin{proof}
  This is the declared model definition of length Value In.
\end{proof}
\begin{definition}[meters Value]
  \label{def:physics:phyx-mini-0046:phyxminiproblems-problemphyxmini0046-metersvalue}
  \lean{PhyXMiniProblems.ProblemPhyXMini0046.metersValue}
  \uses{def:physics:phyx-mini-0046:phyxminiproblems-problemphyxmini0046-dimlength, def:physics:phyx-mini-0046:phyxminiproblems-problemphyxmini0046-lengthvaluein}
  The metre readout of a physical length
\end{definition}
\begin{proof}
  This is the declared model definition of meters Value.
\end{proof}
\begin{definition}[millimeters Value]
  \label{def:physics:phyx-mini-0046:phyxminiproblems-problemphyxmini0046-millimetersvalue}
  \lean{PhyXMiniProblems.ProblemPhyXMini0046.millimetersValue}
  \uses{def:physics:phyx-mini-0046:phyxminiproblems-problemphyxmini0046-dimlength, def:physics:phyx-mini-0046:phyxminiproblems-problemphyxmini0046-lengthvaluein}
  The millimetre readout used for the slit separation and fringe displacement
\end{definition}
\begin{proof}
  This is the declared model definition of millimeters Value.
\end{proof}
\begin{definition}[nanometers Value]
  \label{def:physics:phyx-mini-0046:phyxminiproblems-problemphyxmini0046-nanometersvalue}
  \lean{PhyXMiniProblems.ProblemPhyXMini0046.nanometersValue}
  \uses{def:physics:phyx-mini-0046:phyxminiproblems-problemphyxmini0046-dimlength, def:physics:phyx-mini-0046:phyxminiproblems-problemphyxmini0046-lengthvaluein}
  The nanometre readout used by the answer choices
\end{definition}
\begin{proof}
  This is the declared model definition of nanometers Value.
\end{proof}
\begin{definition}[Double Slit Apparatus]
  \label{def:physics:phyx-mini-0046:phyxminiproblems-problemphyxmini0046-doubleslitapparatus}
  \lean{PhyXMiniProblems.ProblemPhyXMini0046.DoubleSlitApparatus}
  \uses{def:physics:phyx-mini-0046:phyxminiproblems-problemphyxmini0046-dimlength}
  The monochromatic light, the pair of slits, and the screen in the depicted apparatus. screenDistanceAlongX is the figure's horizontal distance R, while slitSeparation is its label d
\end{definition}
\begin{proof}
  This is the declared model definition of Double Slit Apparatus.
\end{proof}
\begin{definition}[Bright Fringe Observation]
  \label{def:physics:phyx-mini-0046:phyxminiproblems-problemphyxmini0046-brightfringeobservation}
  \lean{PhyXMiniProblems.ProblemPhyXMini0046.BrightFringeObservation}
  \uses{def:physics:phyx-mini-0046:phyxminiproblems-problemphyxmini0046-dimlength}
  The signed order, vertical screen coordinate, and ray angle of one bright fringe. The central bright fringe is at vertical coordinate zero and order zero; the observed fringe in the figure lies above it at order +3
\end{definition}
\begin{proof}
  This is the declared model definition of Bright Fringe Observation.
\end{proof}
\begin{definition}[Has Stated Figure Readouts]
  \label{def:physics:phyx-mini-0046:phyxminiproblems-problemphyxmini0046-hasstatedfigurereadouts}
  \lean{PhyXMiniProblems.ProblemPhyXMini0046.HasStatedFigureReadouts}
  \uses{def:physics:phyx-mini-0046:phyxminiproblems-problemphyxmini0046-metersvalue, def:physics:phyx-mini-0046:phyxminiproblems-problemphyxmini0046-millimetersvalue, def:physics:phyx-mini-0046:phyxminiproblems-problemphyxmini0046-doubleslitapparatus, def:physics:phyx-mini-0046:phyxminiproblems-problemphyxmini0046-brightfringeobservation}
  The four numerical labels read from the problem text and primary figure. This predicate deliberately contains no readout of the unknown wavelength
\end{definition}
\begin{proof}
  This is the declared model definition of Has Stated Figure Readouts.
\end{proof}
\begin{definition}[Satisfies Double Slit Bright Fringe Laws]
  \label{def:physics:phyx-mini-0046:phyxminiproblems-problemphyxmini0046-satisfiesdoubleslitbrightfringelaws}
  \lean{PhyXMiniProblems.ProblemPhyXMini0046.SatisfiesDoubleSlitBrightFringeLaws}
  \uses{def:physics:phyx-mini-0046:phyxminiproblems-problemphyxmini0046-metersvalue, def:physics:phyx-mini-0046:phyxminiproblems-problemphyxmini0046-doubleslitapparatus, def:physics:phyx-mini-0046:phyxminiproblems-problemphyxmini0046-brightfringeobservation}
  The exact straight-line screen geometry and constructive double-slit interference law for the observed bright fringe: y = R tan θ, and d sin θ = m λ. The equations use metre readouts so every term in each dimensional equation has the same scalar unit. They contain no numerical value for the requested wavelength
\end{definition}
\begin{proof}
  This is the declared model definition of Satisfies Double Slit Bright Fringe Laws.
\end{proof}
\begin{definition}[Answer Choice]
  \label{def:physics:phyx-mini-0046:phyxminiproblems-problemphyxmini0046-answerchoice}
  \lean{PhyXMiniProblems.ProblemPhyXMini0046.AnswerChoice}
  Labels of the four multiple-choice answers printed with the problem
\end{definition}
\begin{proof}
  This is the declared model definition of Answer Choice.
\end{proof}
\begin{definition}[answer Wavelength Nanometers]
  \label{def:physics:phyx-mini-0046:phyxminiproblems-problemphyxmini0046-answerwavelengthnanometers}
  \lean{PhyXMiniProblems.ProblemPhyXMini0046.answerWavelengthNanometers}
  \uses{def:physics:phyx-mini-0046:phyxminiproblems-problemphyxmini0046-answerchoice}
  Each displayed answer choice, expressed as a whole number of nanometres
\end{definition}
\begin{proof}
  This is the declared model definition of answer Wavelength Nanometers.
\end{proof}
% --- Archon declaration topology end ---
% --- Archon physics formalization source end ---
